\chapter{Conclusions and Future Research}

\section{Conclusions}

This thesis presented RayCastED, a lightweight, fully convolutional, end-to-end model for nuclei detection and segmentation
in histopathology images. Unlike prior methods that rely on dense per-pixel prediction and handcrafted post-processing
(NMS, watershed, or flow-field), RayCastED formulates the task as a sparse regression problem on an anchor grid and produces
final instance masks in a single forward pass. The key findings are summarized below.

\paragraph{Detection and Segmentation Accuracy}
RayCastED achieves competitive performance against state-of-the-art methods on the PanNuke benchmark under a 3-fold
cross-validation protocol. The largest variant (RayCastED-L, 9.44M parameters) attains the fourth-highest overall mPQ
(0.497) and ties for the lowest per-tissue mPQ standard deviation (0.048) among all evaluated methods, indicating the
most consistent panoptic quality across heterogeneous organ types at its parameter tier. This is achieved at $3$--$5\times$
fewer parameters than the baseline cluster (StarDist, CPP-Net, HoVer-NeXt, LSP-DETR) and almost $75\times$ fewer than
LKCell (699.7M parameters), while sharing the same $O(n)$ linear time complexity as LSP-DETR via its
post-hoc-processing-free design. The ablation studies confirm that the Haar wavelet DWT (ResoConv) is the most impactful
architectural component, sweeping all five evaluation metrics against alternative wavelet bases. The ray count ablation
reveals that the model converges to an elliptical default and does not exploit the extra angular precision offered by
denser ray configurations -- a capacity limitation of the small backbone that motivates further research.

\paragraph{Edge Deployment Feasibility}
The smallest variant (RayCastED-S, 0.81M parameters) demonstrates viable inference on the Jetson Orin Nano embedded AI
accelerator. The post-hoc-processing-free design ensures that inference complexity is image-size agnostic, determined
solely by the fixed anchor grid at a given input resolution. This property enables throughput-consistent deployment
across varying tile dimensions, a practical requirement for clinical whole-slide image analysis pipelines.
TensorRT-optimized inference achieves real-time or near-real-time throughput on both desktop and embedded platforms.

\paragraph{Generalization}
Zero-shot evaluation on PUMA, panopTILs, and MoNuSAC reveals reasonable generalization of the PanNuke-trained model to
PUMA (melanoma TILs), but a substantial performance drop on panopTILs and MoNuSAC, indicating that tissue-level and
staining-protocol domain gaps require fine-tuning or retraining for target clinical domains. The RayCast representation
itself is not dataset-specific and can be generated from any polygon annotation format, making the overall pipeline
applicable to diverse data sources.

\section{Future Research}

Several directions warrant further investigation:

\begin{enumerate}
	\item \textbf{Larger backbones to exploit ray density.} The ray-count ablation (Section~\ref{sec:ablation-rays})
	      revealed that RayCastED-B converges to an elliptical default and does not exploit the extra angular precision
	      offered by denser ray configurations, even though the representation itself supports finer boundary detail
	      (Table~\ref{tab:raycast-iou}). Scaling the backbone -- while preserving the lightweight design philosophy --
	      could allow the model to break out of this elliptical prior and discriminate finer boundary features on
	      irregularly shaped nuclei.
	\item \textbf{Adaptive ray density.} The current fixed ray count ($R = 64$) applies uniformly to all nuclei regardless
	      of size. An adaptive scheme that assigns more rays to larger nuclei and fewer to smaller cells could reduce
	      computational cost while maintaining boundary accuracy for large, clinically relevant nuclei.
	\item \textbf{Non-star-convex representations.} The star-convex constraint limits the RayCast representation to simple
	      polygon shapes. Extending the representation to handle concavities and irregular boundaries would broaden
	      applicability to complex nuclear morphologies.
	\item \textbf{Domain adaptation and fine-tuning.} The generalization results indicate that domain-specific retraining
	      is necessary for robust performance on new tissue types and staining protocols. Lightweight fine-tuning
	      strategies, such as adapter layers or prompt-based adaptation, could enable rapid deployment to new clinical
	      settings while preserving the backbone's learned representations.
	\item \textbf{Investigating the Jetson TensorRT accuracy anomaly.} The Jetson Orin Nano TensorRT configurations
	      consistently outperform the desktop RTX 5070 Ti on bPQ and mPQ, an effect not observed on the older RTX 2080
	      Super running the same TensorRT pipeline. Since quantization alone does not account for this discrepancy, a
	      systematic investigation of architecture-specific TensorRT graph optimizations, memory layout differences, and
	      numerical precision paths could reveal optimizations generalizable to other deployments.
	\item \textbf{INT8 quantization.} Further acceleration via INT8 quantization could reduce inference latency and memory
	      footprint on edge devices. Quantization-aware training or calibration-based post-training quantization should be
	      evaluated for impact on detection and segmentation accuracy.
	\item \textbf{Whole-slide image integration.} The tile-based processing pipeline can be extended to end-to-end WSI
	      analysis by coupling RayCastED with a tissue detection front-end and a tile-level result stitching and
	      deduplication back-end. This would enable direct clinical deployment for tasks such as TILs scoring and tumor
	      microenvironment analysis.
	\item \textbf{Extension to other medical tasks.} The RayCast representation and the post-hoc-processing-free detection
	      pipeline are not specific to nuclei. The same framework can be applied to other medical instance-level tasks such
	      as glomeruli detection in renal pathology, gland segmentation in colon histology, and cell counting in
	      cytology.
\end{enumerate}
